\chapter{zmail库}
\href{https://github.com/zhangyunhao116/zmail}{zmail库}不需要手动添加服务器地址、端口以及适合的协议，zmail会帮你完成。此外，使用一个python字典来代表邮件内容也更符合直觉。

\textcolor{red}{\textbf{注意}}:2022年5月7月，zmail库的作者停止维护。

\section{安装zmail库}
\begin{lstlisting}[style=demostyle,numbers=none]
pip install zmail
\end{lstlisting}

\section{使用zmail库}
\subsection{发送邮件}
\begin{lstlisting}[style=demostyle,numbers=none]
import zmail

mail = {
    "subject": "测试邮件", # 邮件主题
    "content": "这是一封测试邮件。", # 邮件正文
    "attachments": ["D:/Practice/python/zmail/test.txt", "D:/Practice/python/zmail/test.py"] # 附件列表
}

password = "087SD4C3B6E3F6G7H8I9J0KLMNOPQRS" # 授权码
server = zmail.server("master@qq.com", password) # 登录邮箱
server.send_mail("slave@qq.com", mail) # 发送邮件
print("邮件发送成功！")
\end{lstlisting}

\subsection{接收邮件}
\begin{lstlisting}[style=demostyle,numbers=none]
import zmail

password = "087SD4C3B6E3F6G7H8I9J0KLMNOPQRS" # 授权码
server = zmail.server("slave@qq.com", password) # 登录邮箱
# mail = server.receive() # 接收邮件
last_mail = server.latest() # 接收最新的一封邮件
print(mail)
\end{lstlisting}
